
\section{Multiset Automata}
\label{Sect_DMA}

Preconditions and conditions in programs are represented by
deterministic multiset automata (DMAs).
A DMA has a finite set of states. 
We use a standard bottom-up semantics (see \cite{TATA2007}),
and we will assume that tree automata are \emph{complete}, i.e.
never get stuck. 
We will merge all preconditions and conditions into a single
DMA. This automaton will use different sets of accepting states
for different (pre)conditions, 
so we do not make the accepting states part of the
automaton. 

On tuples, a DMA behaves as usual: For 
$ ( t_1, \ldots, t_n ), $ once the states for the
$ t_i $ are known, the state for the
tuple can be determined. Since the possible states
for the $ t_i $ are finite in number, and the maximal
value of $ n $ is a priori known, the behaviour
of a DMA on tuples can be finitely expressed.

Arrays will be treated as multisets. In our observation,
all types and subtypes that 
we want to express, can be expressed as set properties
(either at least one element, or all elements must have some property). 
If a type property
cannot be expressed as a set, this means that order matters.
In that case, values at different positions have different meanings,
which implies that they should have been in different fields.
In order to define set transitions, it is sufficient to 
define a state for the empty set, and a binary function for
insertion of a single element. Insertion must be idempotent, 
and order insensitive. This results in two axioms. However, it turned out
that idempotence was not needed, so we omitted it.
This results in a limited form of multiset semantics. 
As a consequence, DMAs can count a bit, but 
with an upperbound, because the set of states is always finite.

\begin{definition}
   \label{Def_DMA}
   We define a \emph{deterministic multiset automaton (DMA)}
   $ {\cal A} $ as a triple
   $ ( Q, \delta, N ), $ where $ Q $ is a finite (nonempty) set of states,
   and $ \delta $ is a transition function with the following
   properties:
   \begin{itemize}
   \item
      If $ c $ is a constant of one of the primitive types
      defined in Definition~\ref{Def_primtype}, then
      $ \delta(c) \in Q. $
   \item
      If $ q_1, \ldots, q_n \in Q $ and $ 0 \leq n \leq N, $
      then
      $ \delta( q_1, \ldots, q_n ) \in Q. $
   \item
      $ \delta( ) \in Q, $
   \item
      If $ q_1,q_2 \in Q, $ then
      $ \delta( q_1; q_2 ) \in Q. $
   \end{itemize}
   The intended meaning of $ ; $ is multiset insertion, with
   the inserted element to the right.
   Technically, $ \delta $ should be viewed as distinct overloads,
   namely
   $ \delta_T $ for primitive type $ T, $ \
   $ \delta_{()} $ for tuples,
   $ \delta_{\emptyset} $ for the empty multiset, and $ \delta_{;} $
   for multiset insertion.
   W.r.t to multiset insertion,
   $ \delta $ must have the following property,
   expressing irrelevance of insertion order: \ 
   $ 
       \delta( \delta( q_1; q_2 ); q_3 ) =
       \delta( \delta( q_1; q_3 ); q_2 ).
   $ 
\end{definition}
As already said before, we are not requiring the property
$ \delta( \delta( q_1, q_2 ); q_2 ) = \delta( q_1; q_2 ), $ which
would enforce set semantics. 

We implicitly assume that $ \delta_T $ on primitive types is
easily representable. The problem is that on an infinite, primitive type,
like $ T = {\bf integer}, \ \ \delta_T $ could represent an 
undecidable property. We assume that on a primitive type $ T, $ 
$ \delta_T $ can be represented by a finite set of values or intervals.

Using a DMA $ {\cal A}, $ every term can be mapped to a state of
$ {\cal A}. $ We will call the resulting function $ Q_{\cal A}: $
\begin{definition}
   For a deterministic tree automaton $ {\cal A} = ( Q, \delta, N ) $,
   we define a function $ Q_{\cal A} $ from data terms to $ Q $
   by recursion as follows:
   \begin{itemize}
   \item
      If $ c $ is a constant of primitive type,
      then $ Q_{\cal A}(c) = \delta(c) $.
   \item
      $ Q_{\cal A}( \ ( d_1, \ldots, d_n ) \ ) =
         \delta( \, Q_{\cal A}(d_1),\ldots,Q_{\cal A}(d_n) \, ) $.
   \item
      $ Q_{\cal A}( \, [ \, ] \, ) = \delta( \, ) $.
   \item
     If $ n > 0, $ then
     $ Q_{\cal A}( \, [ d_1, \ldots, d_n ] \, ) =
          \delta( \, Q_{\cal A}( \, [ d_1, \ldots, d_{n-1} ] \, ); \,
                     Q_{\cal A}(d_n) \, ). $
   \end{itemize}
\end{definition}
Our definition of DMA does not include accepting states.
Our type checking algorithm starts by collecting all possible overloads
for each imprecise name in a translation unit.
Each overload has a type/precondition which
is defined by a DMA. We compute the product of all these automata,
where each contributing automata defines a set of accepting states
in the product automaton. As a consequene, we end up with a single
DMA that is used with different sets of accepting states.

\begin{definition}
   \label{Def_acc_cond}
   Let $ {\cal A} = ( Q, \delta, N ) $ be a DMA.
   An \emph{acceptance condition} $ F $ is a unary predicate on $ Q. $
   An \emph{acceptance relation} $ R $ is a binary predicate on $ Q. $
\end{definition}
Acceptance conditions are used for expressing conditions and
preconditions. Acceptance relations are used for expressing 
conditions only.
\begin{example}
   Propositional formulas are expressed by the following DMA using
   acceptance condition $ \{ \state{form} \}: $ 
   \[
      \begin{array}{lll}
         \delta( \sel{atom} ) & = & \state{atom} \\
         \delta( \sel{neg} ) & = & \state{neg} \\
         \delta( \sel{implies} ) = 
            \delta( \sel{equiv} ) & = & \state{binary} \\
         \delta( \sel{and} ) = 
            \delta( \sel{or} ) & = & \state{nary} \\
         \mbox{ for every {\bf char} } c : \ \ \delta( c ) & = & \state{char} \\
         \\
         \delta( \, ) & = & \state{empty} \\
         \delta( \state{empty}; \state{char} ) = 
         \delta( \state{allchar}; \state{char} ) & = & \state{allchar} \\
         \delta( \state{empty}; \state{form} ) =
         \delta( \state{allform}; \state{form} ) & = & \state{allform} \\
         \\  
         \delta( \state{atom}, \state{allchar} ) & = & \state{prop} \\
         \delta( \state{neg}, \state{prop} ) & = & \state{prop} \\
         \delta( \state{binary}, \state{prop}, \state{prop} ) & = & \state{prop} \\
         \delta( \state{nary}, \state{allprop} ) & = & \state{prop} \\
      \end{array}
   \]
   All transitions that are not listed go to a trap state $ \state{trap}. $ 
\end{example} 
In the example, there is a single trap state $ \state{trap}. $ In practice,
it is convenient to use more than one state for data terms that are not part
of a formula, so that not all information is lost when a formula
is stored in some other data structure.

\noindent
Taking the product of two automata is easy in principle, but
one has to deal with the case where the two automata have different
width:

\begin{definition}
   \label{Def_inflation}
   Let $ {\cal A} = ( Q, \delta, N ). $
   Assume that $ N' > N. $
   The result of \emph{inflating} $ {\cal A} $ to $ N' $
   is the DMA $ {\cal A}' = ( Q', \delta', N' ), $ satisfying: 
   \begin{itemize}
   \item
      $ Q' = \{ Q \} \cup \{ q' \}, $ with $ q' \not \in Q $ a new state.
   \item
      If $ n \leq N, $ then $ \delta'( q_1, \ldots, q_n ) = \delta(q_1, \ldots, q_n). $
      If $ N < n \leq N', $ then $ \delta( q_1, \ldots, q_n ) = q'. $
   \item
      If one of the $ q_i = q', $ then $ \delta(q_1, \ldots, q_n ) = q'. $ 
   \item
      If $ q_1 = q' $ or $ q_2 = q', $ then $ \delta(q_1;q_2) = q'. $ 
   \end{itemize}
\end{definition}

\begin{definition}
   \label{Def_product_automaton}
   Let $ {\cal A}_1 = ( Q_1, \delta_1, N_1) $ and
   $ {\cal A}_2 = ( Q_2, \delta_2, N_2 ) $
   be deterministic tree automata.
   We define the product automaton
   $ {\cal A}_1 \times {\cal A}_2 =
        ( \, Q_1 \times Q_2, \ \delta', \ {\bf max}(N_1,N_2) \, ) \ $
   as follows:
   If $ N_1 \not = N_2, $ then first inflate the smaller
   automata to $ \max(N_1,N_2). $
   After that, construct $ \delta' $ as follows:
   \begin{itemize}
   \item
      For a constant $c$ of primitive type,
      $ \delta'(c) = ( \, \delta_1(c), \delta_2(c) \, ). $
   \item
      $ \delta'( \, (q_1,q'_1), \ldots, (q_n,q'_n) \, ) =
      ( \, \delta_1( q_1, \ldots, q_n), \, \delta_2( q'_1, \ldots, q'_n ) \, ). $
   \item
      $ \delta'( ) = ( \, \delta_1( ), \delta_2( ) \, ). $
   \item
      $ \delta'( \, ( q_1,q'_1); ( q_2, q'_2) \, ) =
        ( \, \delta_1( q_1; q_2), \delta_2( q'_1, q'_2) \, ). $ 
    \end{itemize}
\end{definition}
The product can be straightforwardly generalized to more than two automata.
If $ {\cal A}_1 $ or $ {\cal A}_2 $ has an acceptance condition,
then this acceptance conditions gets 'inherited' by the product automaton
as follows:
\begin{definition}
   \label{Def_acc_proj}
   Let $ F $ be an acceptance condition.
   $ {\bf proj}_i $ is an acceptance condition on tuples
   of arity at least $ i, $ defined from
   $ {\bf proj}_i( \, ( q_1, \ldots, q_n ) \, ) \Leftrightarrow
     F(q_i). $
\end{definition}
Using projection, acceptance conditions can be carried over to the 
product automaton. 
This has one important consequence, namely that
all Boolean combinations of properties definable
by a DMA + acceptance condition can be decided
by constructing the product automata, constructing
a Boolean combination of projections, and verifying it for
all states. At this point, we are relying on the fact that
$ \delta_T $ has low complexity on primitive type $ T. $ 

In particular, we will need to check inclusion between subtypes, 
in cases where more than one overload candidate remains.
In that case, one overload must be strictly more specific than
all other overload candidates.

